\section{Two-sided deficiency for the moving-ceiling experiment}
\label{sec:v24-endpoint-time-deficiency}

The point-process convergence in Section~\ref{sec:v23-full-endpoint-time}
identifies the rare ceiling layer.  For the statistical flagship theorem we
need the stronger statement that this layer is the whole first-order
experiment.  We therefore give parameter-independent Markov kernels in both
directions and a uniform deficiency bound.  This also records explicitly the
effect of the moving-support remainder.

We work first with one design and write \(k=k_n\), \(w=w_\ell\),
\(\rho=\rho_\ell\), \(U=U_\ell\).  On a compact local parameter set \(K\),
write the exact successful density in the reference coordinates as
\[
 q_{n,z}(u,v,r)=\rho_{n,z}(u,v,r)
 \mathbf 1_{\{0<r<w_{n,z}(u,v)\}}.
\]
The fixed-window expansion gives
\begin{equation}\label{eq:v24-deficiency-remainder}
 \varepsilon_{n,K}:=
 \sup_{z\in K}
 \left\{
 k\left\|w_{n,z}-w+k^{-1}Uz\right\|_{C^1}
 +\|\rho_{n,z}-\rho\|_{C^0}
 \right\}\longrightarrow0,
\end{equation}
where the second term may be replaced by its centered first-order smooth
expansion on the bulk; only the displayed convergence rate of the support is
used in the singular layer.  The endpoint collar and all constants below are
uniform on \(K\).

Choose
\[
 R>2+\sup_{z\in K,(u,v)\in\overline D}|U(u,v)z|
\]
and define the \emph{reference layer}
\begin{equation}\label{eq:v24-reference-layer}
 L_{n,R}=\{(u,v,r):|k(w(u,v)-r)|\le R\}.
\end{equation}
The definition contains no unknown parameter.  The fixed face \(r=0\) meets
this layer only when \(0<w\le 2R/k\); as in
\eqref{eq:v23-et-corner-area}, the probability of that corner is
\(O(k^{-2})\) per record.

Let \(p_{n,z}=Q_{n,z}(L_{n,R})\), let \(A_{n,z}\) be the conditional law in
the layer after the scaling \(y=k(w-r)\), and let \(B_{n,z}\) be the
conditional law on \(L_{n,R}^c\).  Let \(\Lambda_z^R\) be the finite measure
on \(D\times[-R,R]\) with density
\begin{equation}\label{eq:v24-Lambda-R}
 d\Lambda_z^R=
ho(u,v)
  \mathbf 1_{\{y>U(u,v)z\}}\mathbf 1_{\{|y|\le R\}}\,du\,dv\,dy.
\end{equation}

\begin{lemma}[Uniform layer and bulk bounds]
\label{lem:v24-layer-bulk-bounds}
Uniformly for \(z\in K\),
\begin{align}
 p_{n,z}&=O(k^{-1}),                                           \label{eq:v24-layer-p}\\
 \left\|\,k\,Q_{n,z}\circ(u,v,r\mapsto u,v,k(w-r))^{-1}
       \big|_{L_{n,R}}-\Lambda_z^R\right\|_{\rm TV}
 &\le C_K(\varepsilon_{n,K}+k^{-1}),                         \label{eq:v24-layer-measure}\\
 H^2(B_{n,z},B_{n,0})&\le C_Kk^{-2}+o_K(k^{-2}).              \label{eq:v24-bulk-H}
\end{align}
In particular the volume, and hence the one-record probability, of the
symmetric difference created by the error term in
\eqref{eq:v24-deficiency-remainder} is at most
\(C_K\varepsilon_{n,K}/k\).
\end{lemma}

\begin{proof}
The layer has thickness \(2R/k\) over an endpoint set of bounded area, which
proves \eqref{eq:v24-layer-p}.  Away from the codimension-two corner, the
change of variable \(y=k(w-r)\) has Jacobian \(k^{-1}\).  The two ceilings in
reference \(y\)-coordinates differ by at most \(\varepsilon_{n,K}\), and the
trace density differs by \(o_K(1)\).  Fubini's theorem bounds the measure of
the symmetric difference by the endpoint area times the ceiling displacement;
the corner contributes \(O(k^{-1})\) after multiplication by \(k\).  This
gives \eqref{eq:v24-layer-measure}.

The choice of \(R\) makes the bulk support common to all \(z\in K\).  On this
common set the normalized density ratio has the uniform smooth expansion
\[
 \frac{dB_{n,z}}{dB_{n,0}}
 =1+k^{-1}a_{n,z}+r_{n,z},
 \qquad
 \sup_{z\in K}\|a_{n,z}\|_\infty\le C_K,
 \qquad
 \sup_{z\in K}\|r_{n,z}\|_\infty=o(k^{-1}).
\]
Normalization removes the first-order mean.  The elementary inequality
\((\sqrt{1+x}-1)^2\le Cx^2\) for small \(x\) then gives
\eqref{eq:v24-bulk-H}.  Notice that this is a statement about the conditional
bulk law, not merely about a marginal likelihood fluctuation.
\end{proof}

Let \(\mathsf E_{n,K}^R=\{Q_{n,z}^{\otimes k}:z\in K\}\).  Let
\(\mathsf P_K^R\) be the Poisson experiment
\(\{\operatorname{PPP}(\Lambda_z^R):z\in K\}\).  A parameter-free ancillary
factor may always be adjoined to either experiment.

\begin{theorem}[Explicit two-sided deficiency bound]
\label{thm:v24-two-sided-deficiency}
For every fixed compact \(K\) and every \(R\) as above,
\begin{equation}\label{eq:v24-LeCam-bound}
 \Delta(\mathsf E_{n,K}^R,\mathsf P_K^R)
 \le C_K\left(
       \varepsilon_{n,K}+k^{-1/2}
             \right)+o_K(1),
\end{equation}
where \(\Delta\) is Le Cam distance.  Both deficiencies are realized by
Markov kernels independent of \(z\).

For finitely many independent design batches the same conclusion holds with
the product Poisson process and the sum of the corresponding right-hand
sides.  In particular the conditional successful endpoint--time experiment
of Theorem~\ref{thm:v23-et-poisson} is asymptotically equivalent, uniformly on
compact local parameter sets, to the asserted Poisson boundary-shift
experiment rather than merely converging to it weakly.
\end{theorem}

\begin{proof}
The forward kernel \(K_n^-\) is deterministic: from the \(k\)-sample retain
exactly the records in the reference layer \(L_{n,R}\), replace their residual
coordinate by \(y=k(w-r)\), and forget the rest.  This operation is
parameter independent.  The retained sample is a binomial point process.
By the classical binomial--Poisson coupling, the total variation distance to
a Poisson process with the same one-record intensity is at most
\(Ckp_{n,z}^2=O(k^{-1})\).  Equation
\eqref{eq:v24-layer-measure} and the standard Poisson coupling for two finite
intensity measures add at most
\(C_K(\varepsilon_{n,K}+k^{-1})\).  This bounds the forward deficiency.

For the reverse direction define a second parameter-independent kernel
\(K_n^+\).  Given a point configuration \(N\) from the Poisson experiment,
let \(m=|N|\).  If \(m\le k\), map its points back to layer records through
\(r=w-y/k\), draw \(k-m\) independent records from the \emph{reference}
bulk law \(B_{n,0}\), and randomly interleave the two collections by a
uniform permutation.  If \(m>k\), output an arbitrary fixed sample.  The
Poisson mean is bounded uniformly on \(K\), so
\(\sup_{z\in K}\Pr_z\{m>k\}=o(k^{-A})\) for every fixed \(A\).

For an iid sample from \(Q_{n,z}\), conditional on the layer count \(m\), the
layer records are iid with law \(A_{n,z}\), the remaining records are iid with
law \(B_{n,z}\), and these two collections are independent before the random
interleaving.  The binomial layer count and layer locations are already
coupled to the Poisson input with error
\(C_K(\varepsilon_{n,K}+k^{-1})\).  For the bulk,
product Hellinger subadditivity and \eqref{eq:v24-bulk-H} give, uniformly in
\(m\),
\[
 H^2(B_{n,z}^{\otimes(k-m)},B_{n,0}^{\otimes(k-m)})
 \le (k-m)H^2(B_{n,z},B_{n,0})=O_K(k^{-1}),
\]
so its total-variation error is \(O_K(k^{-1/2})\).  The random permutation is
a common Markov kernel and cannot increase distance.  Summing these errors
proves \eqref{eq:v24-LeCam-bound}.  Tensorization over the finite list of
design batches proves the product assertion.
\end{proof}

\begin{remark}[Projective Poisson experiment]
\label{rem:v24-projective-poisson}
Let \(\mathcal M_{\rm loc}(D\times\mathbb R)\) carry the usual vague topology.
The Poisson laws with intensity
\(\rho(u,v)\mathbf1_{\{y>U(u,v)z\}}du\,dv\,dy\) form a projective experiment:
restriction from \([-R',R']\) to \([-R,R]\), \(R'>R\), is a
parameter-independent Markov kernel, and the finite-window experiments are
consistent.  For a fixed compact parameter set one finite \(R\) larger than
all possible first-order ceiling displacements contains the singular
information needed above.  Thus the phrase ``increase \(R\)'' refers to this
standard projective system, not to an informal limiting operation.
\end{remark}
